Neural networks have proved themselves an important tool in the analysis of high-dimensional and complex data being able to learn patterns and relationships.
In many physical modelling problems, existing mathematical models exist that characterise the evolution of a system but the models require parameters.
Direct observations of these model parameters often do not exist and must therefore be estimated from data.
For that challenge, neural networks could be constrained to learn the physical laws of a system to estimate parameters for the mathematical models.
These physics-informed neural networks are designed to leverage the prior information concerning a systems evolution to capture the essential patterns of even low amount of data.

However, in many problems, direct measurements from the quantity of interest are not available.
But even then, there may be partial measurements or measurments from one or more related quantities that provide useful information.
In these situtations, Bayesian filtering techniques such as the well-known Kalman filter provide a powerful method to probabilistically estimate the dynamic state of a system in terms of the actual quantity of interest.  

In this thesis, a neural network is constructed to predict advection fields from two-dimensional data, and it is used as a state model operator in an ensemble based Kalman filter.
The work is done as a part of a research project in the Meteorological Research Applications group at the Finnish Meteorological Institute.
Consequently, the model is trained and tested on meteorological satellite data of physical properties of clouds but the general model is designed to be independent of the used data.

In meteorological short-term forecasting, or nowcasting, the so-called physics-free models have been extensively researched in e.g.\ percipitation forecasts.~\cite{rainnet, metnet}
While these methods have shown some efficiency in forecasting applications, the neural networks often lack of interpretability and require large amounts of training data.
As actual physical numerical weather prediction models lack the ability to generate instanteneous high-resolution predictions, the currently used methods are crude mathematical models such as the advection based model studied in this thesis.
As a combination of the physics-free neural networks and the current advection based models, informing the neural network with the advection equation is explored. 
Similar work on physics-informed advection based neural networks have previously been done in~\cite{deepide, debezenac} with sea surface temperature data but in this work the objective is to generate smooth advection fields to better meet the requirements of small-scale cloud data. 
As already tested in~\cite{deepide}, the neural network model is integrated with the Kalman filter model but instead of the goal of probabilistic forecasting, the emphasis is on filling incomplete observations.

The main contributions of this thesis is the repurposing of an existing image recognition neural network architecture to a physics-informed neural network learning to produce smooth advection fields, and the integration of it to a high-dimensional ensemble based Kalman filtering scheme.
The scope of this work is primarily on the exploration and quick validation of new and alternative methods.
Consequently, the neural network development is mostly motivated by general data independent physical interpretability with less focus on performance based tuning.
Moreover, the Kalman filter is set up in an artificial setting in order to provide visually appropriate validation.  
Nevertheless, the neural network model is benchmarked against a well-used classical method used in advective meteorological nowcasting using a test partition of the used training dataset.

The thesis unfolds as follows.
First, some background material on the theory behind the used methods is presented.
Due to the wide range of the used methods and the theory behind them, naturally a complete primer is outside of the scope of this work.
As a compromise based on the spirit of this work, the background section is divided to three parts consisting of characteristics of spatio-temporal data, neural network basics, and Bayesian filtering theory. 
Even though the model was constructed hand in hand with the used dataset, due to it having been designed for general advective spatio-temporal data the proposed model is presented next independently. 
Finally, numerical experiments and verification of the trained model is presented with both the cloud data and synthetic test images separately for the neural network model and the Kalman filter application.
In the very last section, an overall summary of the work is discussed along with possible directions for further research from the basis of this work.
